%-----------------------------------------------------------------------------%
\chapter{\babTiga}
\label{bab:3}
%-----------------------------------------------------------------------------%
Pelaksanaan kerja praktik sebagai Mobile Developer Intern di PT TASPEN (Persero)
memberikan pengalaman berharga dalam mengaplikasikan pengetahuan dan keterampilan
yang didapatkan selama masa perkuliahan ke dalam lingkungan kerja profesional.
Selama periode ini, pelaksana berkesempatan untuk mengembangkan fitur-fitur
penting pada aplikasi internal TASPEN bernama Easy Mobile, yang dirancang untuk
mempermudah kegiatan operasional karyawan. Pelaksanaan kerja praktik ini tidak
hanya mengasah kemampuan teknis pelaksana dalam menggunakan Flutter, tetapi juga
melatih keterampilan non-teknis yang penting, seperti kemampuan beradaptasi,
komunikasi, inisiatif, dan manajemen waktu.

Meskipun dihadapkan pada beberapa tantangan, seperti ketiadaan dokumentasi API
yang lengkap dan kompleksitas \f{business logic} aplikasi, pengalaman ini melatih
pelaksana untuk proaktif dalam mencari solusi dan berkomunikasi secara efektif
dengan mentor. Dengan dukungan tim dan mentor yang baik, pelaksana dapat
menyelesaikan tugas-tugas yang diberikan dan mendapatkan pemahaman yang lebih
dalam mengenai proses pengembangan aplikasi. Secara keseluruhan, pelaksanaan
kerja praktik ini memberikan wawasan yang luas tentang pengembangan perangkat
lunak dan mampu mempersiapkan pelaksana untuk menghadapi tantangan di dunia kerja
yang sesungguhnya.



%-----------------------------------------------------------------------------%
\section{Kesimpulan}
\label{sec:kesimpulan}
%-----------------------------------------------------------------------------%
Pelaksanaan kerja praktik ini memberikan pengalaman berharga bagi pelaksana
dalam mengaplikasikan pengetahuan dan keterampilan yang diperoleh selama masa
perkuliahan ke lingkungan profesional. Selama periode ini, pelaksana
mengembangkan fitur-fitur penting pada aplikasi Easy Mobile yang melatih
kemampuan teknis pelaksana dalam menggunakan Flutter serta keterampilan
nonteknis seperti adaptasi, komunikasi, inisiatif, dan manajemen waktu.
Meskipun menghadapi beberapa tantangan, pelaksana berhasil menyelesaikan
tugas-tugas yang diberikan. Secara keseluruhan, pengalaman ini memberikan
wawasan mendalam tentang pengembangan perangkat lunak dan mempersiapkan
pelaksana untuk menghadapi dunia kerja di masa depan.



%-----------------------------------------------------------------------------%
\section{Saran}
\label{sec:saran}
%-----------------------------------------------------------------------------%
Pada subbab ini akan dijelaskan saran dalam pelaksanaan kerja praktik bagi
tempat kerja praktik, pelaksana kerja praktik berikutnya, dan Fasilkom UI.

%-----------------------------------------------------------------------------%
\subsection{Untuk Tempat Kerja Praktik}
\label{sec:saranTempatKp}
%-----------------------------------------------------------------------------%
Tempat kerja praktik dapat mulai membiasakan tim pengembang perangkat lunaknya
untuk menyusun dokumentasi yang komprehensif pada program yang dibuat agar
proses pengembangan aplikasi ke depannya dapat berlangsung dengan lebih cepat
dan efisien. Hal ini akan sangat membantu para pengembang baru atau peserta
kerja praktik berikutnya untuk memahami struktur dan fungsi API dengan lebih
baik sehingga mereka tidak perlu lagi mengandalkan Postman \f{collection} atau
bertanya langsung kepada mentor untuk setiap detail. Selain itu, menyediakan
panduan atau diagram mengenai \f{user flow} yang kompleks untuk setiap fitur
akan meminimalkan munculnya kebingungan dan memungkinkan pengembang untuk
merancang kode yang lebih terstruktur dan rapi dari awal.

%-----------------------------------------------------------------------------%
\subsection{Untuk Pelaksana Kerja Praktik Berikutnya}
\label{sec:saranPelaksanaBerikutnya}
%-----------------------------------------------------------------------------%
Pelaksana kerja praktik berikutnya disarankan untuk selalu proaktif dan
memiliki inisiatif di lingkungan kerjanya. Lingkungan kerja sering kali terasa
dinamis dan mungkin tidak selalu menyediakan sumber daya yang lengkap untuk
mendukung pekerjaan yang dilakukan. Oleh karena itu, penting untuk tidak hanya
mengandalkan sumber daya yang ada, tetapi juga mengambil inisiatif untuk
mencari solusi. Jika dihadapkan pada keterbatasan, ambil langkah proaktif
dengan berdiskusi langsung dengan mentor atau tim untuk mendapatkan klarifikasi
yang dibutuhkan. Manfaatkan pengalaman dalam bekerja untuk mengasah kemampuan
analitis. Jika menemui alur kerja atau proses bisnis yang kompleks, coba
lakukan analisis mandiri terlebih dahulu. Pengalaman ini dapat melatih
pelaksana selanjutnya dalam memahami logika bisnis yang rumit dan
menerjemahkannya menjadi implementasi yang efisien.

%-----------------------------------------------------------------------------%
\subsection{Untuk Fasilkom UI}
\label{sec:saranFasilkom}
%-----------------------------------------------------------------------------%
Fasilkom UI diharapkan dapat mempertahankan kebijakan serta pelaksanaan dari
mata kuliah Kerja Praktiknya. Pelaksanaan kerja praktik mampu memberikan
pengalaman bagi mahasiswa untuk mengaplikasikan pembelajaran yang sudah
dilakukan semasa kuliah pada lingkungan profesional dengan dampak yang nyata.
Selain itu, kerja praktik juga tidak hanya melatih keterampilan teknis
mahasiswa, tetapi juga mampu melatih keterampilan non teknis yang tidak kalah
penting dalam dunia kerja, seperti komunikasi dan kolaborasi dengan pekerja
lainnya.
